\documentclass[10pt,a4paper]{amsart}
\usepackage[margin=19mm]{geometry}
\usepackage{amsmath,amssymb,mathtools}
\usepackage[scr=rsfs,cal=euler]{mathalfa}
\usepackage{xfrac}
\usepackage[unicode,hidelinks,bookmarks=false]{hyperref}
\newcommand{\ie}{i.e.\ }
\newcommand{\cf}{cf.\ }
\newcommand{\resp}{respectively\ }
\newcommand{\Subsection}{\subsection}
\providecommand{\nobreakdash}{\nobreak\hbox{-}}
\setlength{\emergencystretch}{2em}
\multlinegap0pt
\usepackage{algorithm}\usepackage{algpseudocode}
\usepackage{amssymb}
\usepackage[normalem]{ulem}
\usepackage{xcolor}
\usepackage{mathtools}
\usepackage{float}
\usepackage{tikz}
\newtheorem{proposition}{Proposition}
\newcommand{\RR}[1]{\mathbb{R}^{#1}}
\newcommand{\romDim}{K}
\title{NN-OpInf: an operator inference approach using structure-preserving composable neural networks}
\author{Eric Parish, Anthony Gruber, Patrick Blonigan, Irina Tezaur}
\date{Source: arXiv:2603.08488v1 (CC BY 4.0)}
\begin{document}
\maketitle

\noindent\emph{Source and licence.} NN-OpInf: an operator inference approach using structure-preserving composable neural networks. Version arXiv:2603.08488v1 (CC BY 4.0), distributed by arXiv under the Creative Commons Attribution 4.0 International licence, http://creativecommons.org/licenses/by/4.0/, as recorded in the arXiv licence metadata for this exact version and shown on the abstract page https://arxiv.org/abs/2603.08488v1. Full licence notice: \texttt{licenses/arxiv-2603.08488-CC-BY-4.0.txt}.\par\smallskip

\noindent\textit{Editorial scope.} Counted unit: the article's proposition proposition (source0.tex lines 914--921). The article's complete original proof of it is delivered with it (source0.tex lines 922--924, one \texttt{proof} environment of 3 lines). No mathematics is written, repaired or re-derived: the statement and the proof are the article's own lines, in the article's own notation, and no retained line is substituted or re-flowed.  No further source-local context is needed: every cross-reference of every retained line resolves inside the two delivered spans.  Nothing was omitted from any retained span, so the package carries no omission record.  The retained lines cite 1 works, whose entries are transcribed verbatim from the article's own bibliography source in the e-print and delivered with short editorial labels recorded in the item's reference mapping.  No retained line contains a figure environment, an image-inclusion command or drawing code, so no image is delivered and the package declares no asset dependency.  Editorial formatting only, in addition: an AMS \texttt{amsart} preamble that restates the article's own declarations and macros used by the retained lines, namely the environments \texttt{proposition} on separate counters, together with the standard packages those lines need.  The retained environments print in this document as Proposition 1 in order of appearance, whereas the article prints the same items with the numbering of its own full document.  The complete native source of the article as distributed in the arXiv e-print for this exact version is bundled unchanged as \texttt{sources/196008/source0.tex}.
\begin{proposition}[Convex linear-operator cases]
For the objective $J(\mathbf{A})$, each of the following optimization problems is convex:
\begin{enumerate}
    \item[(i)] $\min_{\mathbf{A}\in\RR{\romDim\times\romDim}} J(\mathbf{A})$ (unconstrained linear operator),
    \item[(ii)] $\min_{\mathbf{S}\in\mathcal{L}_{\mathrm{sl}}} J(\mathbf{S}-\mathbf{S}^{\intercal})$ (skew-symmetric operator with $\mathcal{L}_{\mathrm{sl}}$ the strictly lower-triangular matrices),
    \item[(iii)] $\min_{\mathbf{A}=\mathbf{A}^{\intercal},\,\mathbf{A}\succeq 0} J(\mathbf{A})$ (SPSD operator).
\end{enumerate}
\end{proposition}

\begin{proof}
Case (1) has a convex feasible set and convex objective. In case (2), the map $\mathbf{S}\mapsto \mathbf{S}-\mathbf{S}^{\intercal}$ is linear and $\mathcal{L}_{\mathrm{sl}}$ is a linear subspace, so $J(\mathbf{S}-\mathbf{S}^{\intercal})$ is convex in $\mathbf{S}$. In case (3), the feasible set is the intersection of the symmetric subspace and the PSD cone, hence convex. Therefore, all three problems are convex~\cite{BoydVandenberghe04}.
\end{proof}

\begin{thebibliography}{99}
\bibitem[BoydVa]{BoydVandenberghe04} {\sc S.~Boyd and L.~Vandenberghe}, {\em {Convex} {Optimization}}, Cambridge University Press, 2004.
\end{thebibliography}
\end{document}
